\chapter{บทนำ}

\section{ความเป็นมาและความสำคัญของปัญหา}

ภาคการเกษตรมีบทบาทสำคัญต่อเศรษฐกิจของประเทศไทย ทั้งในด้านการจ้างงานและความมั่นคงทางอาหาร โดยมีแรงงานในภาคการเกษตรมากกว่าร้อยละ 30 ของกำลังแรงงานทั้งหมด อย่างไรก็ตาม ภาคการเกษตรมีสัดส่วนมูลค่าเพิ่มเพียงประมาณร้อยละ 10 ของผลิตภัณฑ์มวลรวมภายในประเทศ และมีแนวโน้มการเติบโตที่ชะลอลงในช่วงหลายทศวรรษที่ผ่านมา \cite{un2020}

เกษตรกรไทยเผชิญกับข้อจำกัดเชิงโครงสร้างหลายประการ ได้แก่ การเปลี่ยนแปลงสภาพภูมิอากาศ การเสื่อมโทรมของทรัพยากร การเข้าถึงเทคโนโลยี และข้อจำกัดด้านองค์ความรู้ ซึ่งส่งผลให้ผลิตภาพต่ำและรายได้ไม่มั่นคง \cite{chadarat2562, rueatee2557}

เพื่อรับมือกับปัญหาดังกล่าว ภาครัฐได้ดำเนินโครงการเกษตรกรอัจฉริยะเพื่อยกระดับศักยภาพของเกษตรกรผ่านการพัฒนาทักษะและการใช้เทคโนโลยี \cite{jansuwan2021} อย่างไรก็ตาม ผลลัพธ์ของโครงการยังมีความแตกต่างระหว่างเกษตรกรแต่ละราย เนื่องจากความแตกต่างด้านพื้นฐาน ความพร้อม และบริบทการผลิต

ในเชิงวิชาการ การวิเคราะห์ทักษะของเกษตรกรจำเป็นต้องคำนึงถึงปัญหาความแตกต่างของกลุ่มตัวอย่างและอคติจากการคัดเลือก \cite{rosenbaum1983} ซึ่งอาจทำให้การประเมินผลคลาดเคลื่อน หากใช้วิธีการเปรียบเทียบแบบทั่วไป

ขณะเดียวกัน เทคโนโลยีปัญญาประดิษฐ์ โดยเฉพาะการเรียนรู้ของเครื่อง มีศักยภาพในการวิเคราะห์ข้อมูลที่มีความซับซ้อนและสามารถอธิบายความแตกต่างของผลลัพธ์ในระดับรายบุคคลได้ \cite{aggarwal2018}

ดังนั้น งานวิจัยนี้จึงมุ่งวิเคราะห์ทักษะความสามารถของเกษตรกรโดยใช้เทคโนโลยีปัญญาประดิษฐ์ ภายใต้กรอบการควบคุมอคติด้วยวิธีทางเศรษฐมิติ เพื่อให้สามารถอธิบายทั้งระดับทักษะโดยรวมและความแตกต่างระหว่างกลุ่มเกษตรกร

\section{วัตถุประสงค์ของการวิจัย}

\begin{enumerate}
\item เพื่อวิเคราะห์ทักษะความสามารถของเกษตรกรโดยใช้เทคโนโลยีปัญญาประดิษฐ์  
\item เพื่อศึกษาปัจจัยที่มีผลต่อการพัฒนาทักษะของเกษตรกร  
\item เพื่อจำแนกกลุ่มเกษตรกรเพื่อสนับสนุนการออกแบบนโยบาย  
\end{enumerate}

\section{ขอบเขตการวิจัย}

ศึกษาจากเกษตรกรที่เข้าร่วมและไม่เข้าร่วมโครงการในช่วงปี พ.ศ. 2565–2566 โดยใช้วิธีการจับคู่เพื่อสร้างข้อมูลที่เหมาะสมสำหรับการวิเคราะห์ด้วยเทคโนโลยีปัญญาประดิษฐ์ \cite{rosenbaum1983}

\section{กรอบแนวคิดในการวิจัย}

กำหนดให้ตัวแปรตามคือระดับทักษะของเกษตรกร และใช้เทคนิคการเรียนรู้ของเครื่องในการวิเคราะห์รูปแบบและจำแนกกลุ่ม โดยใช้วิธีทางเศรษฐมิติเป็นเครื่องมือในการเตรียมข้อมูลให้เหมาะสม \cite{wooldridge2009}

\section{ประโยชน์ที่คาดว่าจะได้รับ}

\begin{enumerate}
\item ได้องค์ความรู้ในการวิเคราะห์ทักษะเกษตรกรด้วยปัญญาประดิษฐ์  
\item ได้แนวทางการออกแบบนโยบายแบบจำเพาะกลุ่ม  
\item สนับสนุนการพัฒนาภาคการเกษตรอย่างมีประสิทธิภาพ  
\end{enumerate}